% !TEX program = tectonic
\documentclass[12pt]{article}
\newif\ifanonymous\anonymousfalse
\usepackage{fontspec}
\setmainfont{Times New Roman}
\usepackage{amsmath}
\usepackage{newtxmath}
\usepackage[a4paper,margin=25mm]{geometry}
\usepackage{microtype}
\usepackage{graphicx}
\usepackage{booktabs,tabularx,makecell,array,longtable}
\usepackage{setspace}
\newif\ifsupplementary\supplementaryfalse
\usepackage[authoryear,round]{natbib}
\bibpunct{(}{)}{,}{a}{}{,}
\usepackage{xurl}
\usepackage[hidelinks]{hyperref}
\hypersetup{pdfauthor={},pdftitle={Classifier selection, score estimation and event coverage in inertial pavement screening}}
\usepackage{caption}
\captionsetup{font=small,labelfont=bf,skip=6pt}
\usepackage{enumitem}
\setlist{nosep,leftmargin=1.5em}
\usepackage{titlesec}
\titleformat{\section}{\normalfont\large\bfseries}{\thesection.}{.6em}{}
\titleformat{\subsection}{\normalfont\normalsize\bfseries\itshape}{\thesubsection.}{.6em}{}
\usepackage[switch]{lineno}
\renewcommand\linenumberfont{\normalfont\tiny}
\newcolumntype{Y}{>{\raggedleft\arraybackslash}X}
\newcolumntype{L}[1]{>{\raggedright\arraybackslash}p{#1}}
\newcolumntype{R}[1]{>{\raggedleft\arraybackslash}p{#1}}
\setlength{\parindent}{1.5em}
\setlength{\parskip}{0pt}
\setlength{\emergencystretch}{2em}
\setlength{\tabcolsep}{3.5pt}
\renewcommand{\arraystretch}{1.16}
\clubpenalty=10000\widowpenalty=10000\displaywidowpenalty=10000
\raggedbottom
\makeatletter\setlength{\@fptop}{0pt}\makeatother
\doublespacing
\newcommand{\articletitle}{Classifier selection, score estimation and event coverage in inertial pavement screening}
\newcommand{\authorblock}{Mingyang Li\textsuperscript{1}, Jiabei Liu\textsuperscript{2,*}, Jiayue Gao\textsuperscript{2}\\[6pt]
\textsuperscript{1}Beijing Institute of Technology, Beijing, China\\
\textsuperscript{2}Jinan University--University of Birmingham Joint Institute, Jinan University, Guangzhou, China\\[6pt]
\textsuperscript{*}Corresponding author: Jiabei Liu; \href{mailto:jxl2043@student.bham.ac.uk}{jxl2043@student.bham.ac.uk}}

\begin{document}\singlespacing
\noindent 3 November 2026\par\bigskip
\noindent Dear Editors,\par\medskip
Please consider our research article, ``\articletitle'', for the \emph{International Journal of Pavement Engineering}. The paper addresses evaluation of low-cost sensor-based road-surface screening, with a focus on selecting classifiers for new acquisition conditions and communicating their expected performance.
\par\medskip
The study combines a version-specific audit of 101 RoadSens recordings with nested selection experiments on nine PVS acquisitions. In RoadSens, random-forest macro-F1 changes from 0.9945 under random-row validation to 0.4289 under recording holdout, while temporal inertial summaries improve held-out classification. In PVS, domain-matched validation selects stronger scenario-held-out systems but does not satisfy the declared criterion for performance-estimation accuracy. Disjoint-window and common-budget diagnostics show why these two evaluation objectives should be reported separately. A two-route Mexico/Puebla task provides a contrasting result: the selectors choose identical systems and have small estimation errors. An independent public Karlsruhe source adds a final GPS-date test with spatially buffered development data and matched inner training counts. On that test, date selection improves classification and estimation, while both chosen systems detect seven of 12 annotated pothole intervals. Event and alarm diagnostics connect the evaluation results to screening workload.
\par\medskip
The connection to pavement engineering is the evaluation evidence supporting a screening model's use across survey conditions. The article relates acquisition grouping and label semantics to the distinction between pavement distress, intentional road features and non-event observations, and quantifies the resulting interval coverage and notification burden. It complements work on smartphone-based pavement measurement by examining classification-based screening and the accuracy of validation estimates.
\par\medskip
The submission materials include an anonymous manuscript, separate author details, supplementary results, editable tables and vector figures. Code, protocols and compact results are publicly available. The study received no specific funding, and the authors report no competing interests.
\par\medskip
Thank you for considering this manuscript.\par\bigskip
\noindent Jiabei Liu\par
Corresponding author\par
Jinan University--University of Birmingham Joint Institute, Jinan University\par
\href{mailto:jxl2043@student.bham.ac.uk}{jxl2043@student.bham.ac.uk}
\end{document}
